% !PRINTED 692
\weberpage{692}

Ist umgekehrt für ein Funktionensystem $\eta_1$, $\eta_2$, \dots, $\eta_n$ die
Bedingung
\weq{10}{S(\eta_r\omega_s) = (r,\ s)}
befriedigt, so ist $\eta_r = \varepsilon_r$; denn setzt man
\[ \eta_r = \sideset{^i}{}\sum x_{i,\,r}\varepsilon_i, \]
so folgt aus (9) und (10) $x_{r,\,s} = (r,\ s)$. Demnach sind die Be\-dingungen
(9) und (6) vollständig gleichbedeutend.

Die Größen $\varepsilon_i$ bilden eine Basis von $\Omega$, sind aber nicht
ganze Funktionen. Die Basen (5) und (2) heißen \gesperrt{zueinander
komplementär}.

Die Größen
\weq{11}{\Delta\varepsilon_1,\ \Delta\varepsilon_2,\ \dots,\ \Delta\varepsilon_n}
sind ganze Funktionen, und wenn $t_1$, $t_2$, \dots, $t_n$ Funktionalvariablen
sind, so ist
\weq{12}{\Delta(\varepsilon_1t_1 + \varepsilon_2t_2 + \cdots + \varepsilon_nt_n) = \Delta\varepsilon}
der größte gemeinschaftliche Teiler der Funktionen (11).

Ist
\weq{13}{r = (z - c_1)(z - c_2)(z - c_3)\dots}
das Produkt aller voneinander verschiedenen Linearfaktoren von
$\Delta$ und
\weq{14}{\varkappa = \pi_1\pi_2\pi_3\dots}
das Produkt der verschieden in $\Delta$ aufgehenden Primfunktionale,
so können wir jede durch $r$ teilbare ganze Funktion $\varrho$ von $z$
in $\Omega$ in die Form setzen:
\[ \varrho = r(x_1\varepsilon_1 + x_2\varepsilon_2 + \cdots + x_n\varepsilon_n), \]
worin die $x_1$, $x_2$, \dots, $x_n$ rationale Funktionen von $z$ sind. Dann
erhält man nach (9)
\[ rx_i = S(\varrho\omega_i), \]
und da $\varrho\omega_i$ durch jeden Primfaktor von $r$ teilbar ist, so ist
$S(\varrho\omega_i)$ nach dem Satz § 184, 7. eine durch $r$ teilbare ganze
Funktion in $Z$,
und folglich $x_i$ eine ganze rationale Funktion.

Hieraus folgt, daß
\[ \Delta\varrho = r(x_1\Delta\varepsilon_1 + x_2\Delta\varepsilon_2 + \cdots + x_n\Delta\varepsilon_n) \]
durch $r\varepsilon\Delta$ teilbar ist, und daß sonach $r\varepsilon$ ein in $\varrho$ aufgehendes
ganzes Funktional ist.

Wäre $\pi$ ein nicht in $\varkappa$ enthaltener Primfaktor von $r\varepsilon$, so
könnte man $\varrho$ so annehmen, daß es nicht durch $\pi$ teilbar ist,
%CONT>>
